\chapter{Conclusion and Future Scope}

The development of speech recognition systems for dysarthric speech highlights the gap between current AI capabilities and real-world accessibility needs. Addressing this gap requires not only improvements in acoustic modeling but also the integration of contextual intelligence. This chapter summarizes the outcomes of the proposed system, evaluates its effectiveness, and outlines possible directions for future improvements.

\section{Conclusion}

The project focused on improving speech recognition for dysarthric speech by addressing key challenges such as limited datasets, high acoustic variability, and noise sensitivity. The objectives included generating synthetic dysarthric speech, enhancing robustness through augmentation, fine-tuning a Whisper-based ASR model, and incorporating a semantic correction layer using language models and a Personal Knowledge Graph.

These objectives were implemented through a structured pipeline. Approximately 16,000 synthetic dysarthric speech samples were generated using Text-to-Speech (TTS) and augmentation techniques, resulting in a dataset size of nearly 3.2 GB. A pre-trained Whisper-small model was fine-tuned using parameter-efficient LoRA techniques to adapt to dysarthric speech patterns. Training was performed using NVIDIA Tesla T4 GPUs on Kaggle and Google Colab platforms, with total training sessions spanning approximately 9 hours on Kaggle (dual T4 GPUs) and 5 additional hours on Google Colab.

The results demonstrate a significant improvement in transcription performance. The baseline system achieved a Word Error Rate (WER) of 0.3215, while the fine-tuned model reduced the WER to 0.1612. Compared to the reference baseline paper with a WER of 0.21, the proposed approach achieved further improvement through synthetic augmentation and semantic correction. This corresponds to an approximate improvement of 49.86\% over the initial baseline system and nearly 23.24\% improvement over the reference approach. These results validate the effectiveness of combining acoustic adaptation with contextual reasoning for dysarthric speech recognition.

% TODO: Edit CER values here
% CER Before:
% CER After:

\section{Future Scope}

Despite the improvements achieved, several limitations remain. The availability of large-scale, diverse dysarthric speech datasets is still a major constraint. Synthetic data, while useful, does not fully capture real-world variability across speakers and conditions. Additionally, the semantic correction layer depends on the quality of ASR output and may introduce latency in real-time applications.

Future work can focus on expanding dataset diversity through real-world data collection and improving synthetic speech generation using advanced GAN-based models. Personalization techniques can be explored to adapt models to individual speakers more effectively. Integration with edge devices and Neural Processing Units (NPUs) can enable real-time, privacy-preserving deployment. Further research can also enhance the Knowledge Graph and semantic reasoning capabilities to improve contextual understanding.

\section{Learning Outcomes of the Project}

\begin{itemize}
\item Understanding of speech signal processing and challenges in dysarthric speech recognition
\item Hands-on experience with deep learning frameworks such as PyTorch and Hugging Face Transformers
\item Implementation of parameter-efficient fine-tuning techniques such as LoRA
\item Knowledge of data augmentation and synthetic speech generation methods
\item Experience in evaluating ASR systems using metrics such as Word Error Rate (WER) and Character Error Rate (CER)
\item Understanding the role of language models in semantic correction and contextual reasoning
\item Experience in training and fine-tuning transformer-based ASR models using cloud GPU platforms
\end{itemize}